\documentclass[10pt,a4paper]{amsart}
\usepackage[margin=19mm]{geometry}
\usepackage{amsmath,amssymb,mathtools}
\usepackage[scr=rsfs,cal=euler]{mathalfa}
\usepackage{xfrac}
\usepackage[unicode,hidelinks,bookmarks=false]{hyperref}
\newcommand{\ie}{i.e.\ }
\newcommand{\cf}{cf.\ }
\newcommand{\resp}{respectively\ }
\newcommand{\Subsection}{\subsection}
\providecommand{\nobreakdash}{\nobreak\hbox{-}}
\setlength{\emergencystretch}{2em}
\multlinegap0pt
\usepackage{amssymb}
\usepackage{bbm}
\usepackage{booktabs}
\usepackage{float}
\usepackage{xcolor}
\usepackage{tikz}
\newtheorem{theorem}{Theorem}
\newtheorem{lemma}{Lemma}
\title{$L$-derivatives of the fixed elliptic curve over rank-one imaginary quadratic fields}
\author{Shenghao Hua}
\date{Source: arXiv:2507.20297v3 (CC BY 4.0)}
\begin{document}
\maketitle

\noindent\emph{Source and licence.} $L$-derivatives of the fixed elliptic curve over rank-one imaginary quadratic fields. Version arXiv:2507.20297v3 (CC BY 4.0), distributed by arXiv under the Creative Commons Attribution 4.0 International licence, http://creativecommons.org/licenses/by/4.0/, as recorded in the arXiv licence metadata for this exact version and shown on the abstract page https://arxiv.org/abs/2507.20297v3. Full licence notice: \texttt{licenses/arxiv-2507.20297-CC-BY-4.0.txt}.\par\smallskip

\noindent\textit{Editorial scope.} Counted unit: the article's theorem thm:lowerbound-L (source0.tex lines 277--295). The article's complete original proof of it is delivered with it (source0.tex lines 427--510, one \texttt{proof} environment of 84 lines, which the article places away from the statement). No mathematics is written, repaired or re-derived: the statement and the proof are the article's own lines, in the article's own notation, and no retained line is substituted or re-flowed.  Retained uncounted as the source-local context the proof needs: lines 246--248; lines 305--323.  Nothing was omitted from any retained span, so the package carries no omission record.  The retained lines cite 2 works, whose entries are transcribed verbatim from the article's own bibliography source in the e-print and delivered with short editorial labels recorded in the item's reference mapping.  No retained line contains a figure environment, an image-inclusion command or drawing code, so no image is delivered and the package declares no asset dependency.  Editorial formatting only, in addition: an AMS \texttt{amsart} preamble that restates the article's own declarations and macros used by the retained lines, namely the environments \texttt{theorem}, \texttt{lemma} on separate counters, together with the standard packages those lines need.  The retained environments print in this document as Theorem 1, Lemma 1 in order of appearance, whereas the article prints the same items with the numbering of its own full document.  The complete native source of the article as distributed in the arXiv e-print for this exact version is bundled unchanged as \texttt{sources/182275/source0.tex}.
\begin{equation}\label{eqn:L'EK}
  L'(1, E/\mathbb{Q}(\sqrt{d})) = L(1, E) L'(1, E^{(d)}) + L'(1, E) L(1, E^{(d)}),
\end{equation}

\begin{lemma}\label{lemma:twist}
Let \( u \) be a positive integer such that \( (u\ell, N_0) = 1 \).
For \( w = 1 \), define
\begin{equation*}
  S'(X; u) = \sum_{\substack{d \in \Omega_{-1}(a)}}
  L'\left(1, E^{(d)}\right) \chi_d(u)
  \Phi\left(\frac{-d}{X}\right).
\end{equation*}

Write \( u = u_1 u_2^2 \leq p \) with \( u_1 \) square-free.
Then there exist absolute constants \( C = C(E) \), \( C_1 = C_1(E, \Phi) \), and constants \( C_2(p) \ll 1 \) for all \( p \), such that for any \( \varepsilon > 0 \), we have
\begin{equation*}
  S'(X; u) = \frac{a(u_1)}{4\pi^2 u_1 N_0} \,
  \check{\Phi}(0) X \left( \log \frac{X}{u_1} + C_1 + \sum_{p \mid u} \frac{C_2(p)}{p} \log p \right)
  + O\left(X^{\frac{3}{4} + \varepsilon} u^{1/2}\right).
\end{equation*}

Here, we write \( a(p) = \alpha_p + \beta_p \), where \( \alpha_p \beta_p = p \), and \( |\alpha_p| = |\beta_p| = \sqrt{p} \) for all primes \( p \nmid N \).
\end{lemma}

\begin{theorem}\label{thm:lowerbound-L}
Let \( E \) be a non-CM elliptic curve over \( \mathbb{Q} \) with conductor \( N \) and root number \( w \), and the vanishing order of the \(L(s,E)\) at the central point is at most 1.
Fix a residue class \( a \bmod N_0 \) be a residue class such that \( a \) is a quadratic residue modulo \( N \), \( a \equiv 1 \) or \( 5 \pmod{8} \),
and for any fundamental discriminant $d<0$ with $d\equiv a \pmod{N_0}$, $ w\chi_d(-N)=\epsilon$.
For \( -3>d\in  \Omega_{-w}(a)\), there exist positive constants \( c_1', c_2', c_3' \) such that for any small \( \varepsilon > 0 \) and sufficiently large \( X \), we have
\[
\left|\left\{
 d \in \Omega_{1}(a) \mid 20 < -d \leq X,\
L'(1, E/\mathbb{Q}(\sqrt{d})) \geq c_1'
\right\}\right| \gg X^{1 - \varepsilon},
\]
\[
\left|\left\{
 d \in \Omega_{-1}(a) \mid 20 < -d \leq X,\
L'(1, E/\mathbb{Q}(\sqrt{d}))\geq c_2' \log|d| + c_3'
\right\}\right| \gg X^{1 - \varepsilon}.
\]
Moreover, when \(L(1,E)\neq 0\), this result also applies to \(L'(1, E^{(d)})\).
\end{theorem}

\begin{proof}[Proof of Theorem~\ref{thm:lowerbound-L}]
From \eqref{eqn:L'EK}, when \( w = -1 \),
applying \cite[Proposition~2]{RadziwillSoundararajan2015}, there exists a constant \( D_1 > 0 \), depending on \( E \), such that
\begin{multline*}
\sum_{d \in \Omega_{1}(a)}
L'(1, E/\mathbb{Q}(\sqrt{d}))
\Phi\left(\frac{-d}{X}\right)
=
L'(1, E)
\sum_{d \in \Omega_{1}(a)} L(1, E^{(d)})
\Phi\left(\frac{-d}{X}\right)
\\= D_1 X + O\left(X^{\frac{7}{8}+\varepsilon}\right).
\end{multline*}

When \( w = 1 \), using Lemma~\ref{lemma:twist}, there exist constants \( D_2, D_3 > 0 \), depending on \( E \), such that
\begin{multline*}
\sum_{d \in \Omega_{-1}(a)}
L'(1, E/\mathbb{Q}(\sqrt{d}))
\Phi\left(\frac{-d}{X}\right)
=
L(1, E)
\sum_{d \in \Omega_{1}(a)}
L'(1, E^{(d)})
\Phi\left(\frac{-d}{X}\right)
\\= D_2 X \log X + D_3 X + O\left(X^{\frac{3}{4}+\varepsilon}\right).
\end{multline*}

Moreover, by Heath-Brown's quadratic large sieve~\cite{HeathBrown1995}, in both cases we have
\[
\sum_{d \in \Omega_{-w}(a)}
L'(1, E/\mathbb{Q}(\sqrt{d}))^2
\Phi\left(\frac{-d}{X}\right)
\ll X^{1 + \frac{\varepsilon}{2}}.
\]

Let \( c_i'< \frac{D_i}{2} \) for \( i = 1, 2, 3 \). Then for \( w = -1 \), we have
\[
\sum_{\substack{d \in \Omega_1(a) \\ L'(1, E/\mathbb{Q}(\sqrt{d})) < c_1'}}
L'(1, E/\mathbb{Q}(\sqrt{d}))
\Phi\left(\frac{-d}{X}\right)
< 2c_1' X,
\]
and thus,
\[
\sum_{\substack{d \in \Omega_1(a) \\ L'(1, E/\mathbb{Q}(\sqrt{d})) \geq c_1'}}
L'(1, E/\mathbb{Q}(\sqrt{d}))
\Phi\left(\frac{-d}{X}\right)
> (D_1 - 2c_1') X.
\]

A similar argument applies for \( w = 1 \).

Now let \( k = \frac{\varepsilon}{3} \). By H\"older's inequality, we obtain
\begin{multline*}
\sum_{\substack{d \in \Omega_1(a) \\ L'(1, E/\mathbb{Q}(\sqrt{d})) \geq c_1'}}
\left|L'(1, E/\mathbb{Q}(\sqrt{d}))
\Phi\left(\frac{-d}{X}\right)\right|^k
\\ \gg
\left(
\sum_{\substack{d \in \Omega_1(a) \\ L'(1, E/\mathbb{Q}(\sqrt{d})) \geq c_1'}}
\left|L'(1, E/\mathbb{Q}(\sqrt{d}))
\Phi\left(\frac{-d}{X}\right)\right|
\right)^{2 - k}
\\ \times \left(
\sum_{\substack{d \in \Omega_1(a) \\ L'(1, E/\mathbb{Q}(\sqrt{d})) \geq c_1'}}
\left|L'(1, E/\mathbb{Q}(\sqrt{d}))
\Phi\left(\frac{-d}{X}\right)\right|^2
\right)^{k - 1}
\gg X^{1 - \frac{\varepsilon}{2}}.
\end{multline*}

Using the convexity bound
\[
L(1, E^{(d)}),~ L'(1, E^{(d)})
\ll d^{\frac{1}{2} + \frac{\varepsilon}{2}},
\]
we deduce
\[
\sum_{\substack{d \in \Omega_1(a) \\ L'(1, E/\mathbb{Q}(\sqrt{d})) \geq c_1'}} 1
\gg X^{1 - \varepsilon}.
\]

A similar conclusion holds for \( w = 1 \). This completes the proof.
\end{proof}

\begin{thebibliography}{99}
\bibitem[HeathB]{HeathBrown1995} Heath-Brown, D. R. A mean value estimate for real character sums. \emph{Acta Arith.} 72 (1995), no. 3, 235--275. %
\bibitem[Radziw]{RadziwillSoundararajan2015} Radziwi\l\l, Maksym; Soundararajan, Kannan. Moments and distribution of central $L$-values of quadratic twists of elliptic curves. \emph{Invent. Math.} 202 (2015), no. 3, 1029--1068.
\end{thebibliography}
\end{document}
